\documentclass[~/PhyAI/phyai.tex]{subfiles}
\begin{document}
\chapter{The First Law of Thermodynamics}\label{chap:the_first_law_of_thermodynamics} % (fold)
\section{The Microscopic Model}\label{sec:the_microscopic_model} % (fold)
As you can see here, we are moving further than the zeroth law of thermodynamics. In the previous chapter, we neither know what is inside the substances, nor the essence of heat. Thinking like a physicist, we need to build a physical model to explain many phenomena in nature.
\begin{center}
	\textit{Substance is just the collection of many very tiny "imaginary" particles.}
\end{center}
\begin{figure}[H]
	\begin{center}
		\includegraphics[width=0.6\textwidth]{~/PhyAI/chapter1/1_0.png}
	\end{center}
	\caption{Our hypothesis model of gas and crystal}\label{fig:1.1}
\end{figure}
But why do I call these particles "imaginary"? Because in the $19$th century, scientists had no way of directly observing them. Physicists living that era built the entire foundation of thermodynamics on atomic theory, which may not have convinced everyone until the quantum age. In the early $20$th century, the existence of atoms was finally proven and widely accepted, supported by numerous undeniable physical and mathematical proofs. Finally, in $1955$, the first photograph of atom was taken by Erwin W.Muller and his student.
\begin{theorem}[Atomic Theory]
	\label{Atomic theory}
	Atomic theory is the specific theory that matter is composed of particles called \textbf{atoms}.
\end{theorem}
I think Theorem ~\ref{Atomic theory} should be called axiom, in mathematical language. In the following chapters, we will focus in depth on how this theory can be applied to build the foundations of thermodynamics, rather than proving it correctness.
% section The Microscopic Model (end)
\section{Volume and Pressure}\label{sec:volume_and_pressure}
I think everyone knows what volume is, but it is not so easy to formalize our sense to a definition. Mathematically, we can state that:
\begin{definition}[The First Definition of Volume]
	\label{The First Definition of Volume}
	Volume is a measure of regions in three-dimensional space.
\end{definition}
But I think we can extend the definition above with atomic theory:
\begin{definition}[The Second Definition of Volume]
	\label{The Second Definition of Volume}
	Volume is a measure of the amount of space occupied by many small particles of \textbf{the same type} at equilibrium state.
\end{definition}
\begin{figure}[H]
	\begin{center}
		\includegraphics[width=0.7\textwidth]{~/PhyAI/chapter1/1_1.png}
	\end{center}
	\caption{Definition of Volume}\label{fig:1.2}
\end{figure}
Now you might wonder if you mix $2$ types of completely different gases, such as $\ce{H2}$ and $\ce{He}$, are they the same type? The answer is yes, but until they reach the \textbf{equilibrium state}.
\begin{figure}[H]
	\begin{center}
		\includegraphics[width=0.7\textwidth]{~/PhyAI/chapter1/1_2.png}
	\end{center}
	\caption{Mixing $\ce{H2}$ and $\ce{He}$ together}\label{fig:1.3}
\end{figure}
Reaching the equilibrium state means the density of gas mixture is the same at all locations inside the box.\footnote[1]{This is just my claim.}\footnote[2]{Mixing gas process can release an amount of heat energy. According to the zeroth law of thermodynamics, equilibrium state is reached when the temperature of the mixture is stable.}
\newpage
\indent Unlike volume, \textbf{pressure} is very easy concept to define:
\begin{definition}[Definition of Pressure]
	Pressure is the force applied \textbf{perpendicular} to the surface of an object per unit area over which that force is distributed:$$ P=\frac{F }{A}\text{ (N/m$^2$), or (Pa)} $$
	\label{Definition of Pressure}
\end{definition}
One classic example of pressure is pressing your hand against a small piston filled with gas. If you press hard and fast enough, you can observe small sparks inside the piston. We will be back to explain the phenomenon later.
\begin{figure}[H]
	\begin{center}
		\includegraphics[width=0.6\textwidth]{~/PhyAI/chapter1/1_3.png}
	\end{center}
	\caption{Pressing a small piston}\label{fig:1.4}
\end{figure}
After understanding what volume and pressure are, now it is time to answer our previous question from Chapter ~\ref{chap:the_zeroth_law_of_thermodynamics}, what is Kelvin scale? Imagine you have an elastic balloon, which can be easily stretched. Of course, this balloon is filled fully with gas.
\begin{figure}[H]
	\begin{center}
		\includegraphics[width=0.8\textwidth]{~/PhyAI/chapter1/1_4.png}
	\end{center}
	\caption{For every degree decrease, the volume $V/273$ decreases correspondingly}\label{fig:1.5}
\end{figure}
But because the volume quantity can not go below zero, if the temperature reaches $(T-273)^{\circ}$C, and continues to decrease, the volume will remain at zero. Alternatively, if you hold the volume of the gas fixed, then its pressure will approach zero as the temperature approaches $(T-273)^{\circ}$C. The special value $-273^{\circ}$C is called \textbf{absolute zero}\footnote[1]{This is the basic idea of how absolute zero was discovered.}, and defines the initial point of the \textbf{absolute temperature scale}, named after Lord Kelvin\footnote[2]{How about holding both the volume and pressure of the gas fixed, and decrease temperature?}.
\begin{definition}[Definition of Kelvin Scale]
	\label{Definition of Kelvin Scale}
	The Kelvin scale is an absolute temperature scale that starts at the lowest possible temperature, taken to be $0$K:
	\begin{equation*}
		T_{K}=T_{C}+273^{\circ}C\text{ (K)}
	\end{equation*}
\end{definition}
% section Volume and Pressure (end)
\section{The Ideal Gas Law}\label{sec:the_ideal_gas_law} % (fold)
Discovered by experiments, many properties of a low-density gas can be summarized in the famous idea gas law, where $P$ (N/m$^2$), $V$ (m$^3$), $n$ (mole), $R$ is a universal constant, and $T$ (K).
\begin{theorem}[The Ideal Gas Law]
	\label{The Ideal Gas Law}
	The empirical ideal gas law for \textbf{low-density gas} is:
	$$ PV=nRT $$
\end{theorem}
This section does not have much to discuss beyond introducing an important law of gases and some new constants. As you can see here, the unit of constant R is:
\begin{equation*}
	R=8.31\text{ (J/mol.K)}
\end{equation*}
If we define $N_{A}=6.023.10^{23}$ (mol$^{-1}$) is Avogadro's number, then the law can be rewritten as:
\begin{equation*}
	PV=nRT=\frac{N }{N_{A}}RT=N\mathbf{k_{B}}T
\end{equation*}
The constant $k_{B}$ is called \textbf{Boltzmann's constant}, it plays vital role in thermodynamics:
\begin{equation*}
	k_{B}=\frac{R }{N_{A}}=1.381.10^{-23}\text{ (J/K)}
\end{equation*}
% section The Ideal Gas Law (end)
\section{The Kinetic Theory of Gases}\label{sec:the_kinetic_theory_of_gases} % (fold)
In a gas-filled piston, you already know that gas (or any substance) is simply a collection of tiny, chaotic, uncontrollable particles. If you heat the piston, these particles appear to move faster. In this section, we will use the gas model to quantify the relationship between piston temperature and the velocity of the particles.
\begin{figure}[H]
	\begin{center}
		\includegraphics[width=0.5\textwidth]{~/PhyAI/chapter1/1_5.png}
	\end{center}
	\caption{An oversimplified model of an ideal gas}\label{fig:1_5}
\end{figure}
Since the system of $10^{23}$ particles is impossible to process in the first place, I have to make some assumptions here to simplify our physical model:
\begin{itemize}
	\item My oversimplified model of an ideal gas contains only $1$ particle.
	\item Its speed never changes.
	\item Piston is perfectly smooth, so collisions are always \textbf{elastic}.
\end{itemize}
Directly applying Definition ~\ref{Definition of Pressure} yields:
\begin{equation*}
	\begin{split}
		\overline{P}_{\text{x, on piston}}&=\frac{\overline{F}_{\text{x, on piston}} }{A}=-\frac{\overline{F}_{\text{x, on particle}} }{A}\text{ (Newton's Third Law)}\\
		&=-\frac{m \displaystyle\frac{\Delta v}{\Delta t}}{A}=-\frac{m \displaystyle\frac{(-v_{x})-v_{x}}{2L/v_{x}}}{A}=\frac{mv_{x}^{2}}{V}\\
	\end{split}
\end{equation*}
If our piston contains $N$ particles, then the total pressure on piston by $x$-axis is:
$$ P_{x}=\sum\overline{P}_{\text{x, on piston}}=\frac{\sum m v_{x,i}^{2} }{V}=N \frac{m\overline{v_{x}^{2}} }{V} \Rightarrow P_{x}V=Nm\overline{v_{x}^{2}}$$
Because $N$ is extremely large, and the gas is homogeneous:
$$ P=P_{x}=P_{y}=P_{z}\Rightarrow PV=Nm\overline{v_{x}^{2}} $$
The average velocity of the particles in the $x, y$ and $z$ directions is the same, we obtain
$$ \frac{m\overline{v_{x}^{2}}}{2}=\frac{m\overline{v_{y}^{2}}}{2}=\frac{m\overline{v_{z}^{2}}}{2}=\frac{kT} {2} $$
The average velocity of $N$ particles in three-dimensional space is defined as:
$$ \overline{v^{2}}=\overline{v_{x}^{2}+v_{y}^{2}+v_{z}^{2}}=\frac{v_{x}^{2}+v_{y}^{2}+v_{z}^{2}}{N}=\overline{v_{x}^{2}}+\overline{v_{y}^{2}}+\overline{v_{z}^{2}} $$
The average translational kinetic energy is then:
$$ \overline{K}=\frac{1 }{2}m\overline{v^{2}}=\frac{1 }{2}m\overline{v_{x}^{2}}+\frac{1 }{2}m\overline{v_{y}^{2}}+\frac{1 }{2}m\overline{v_{z}^{2}}=\frac{3kT}{2} $$
We can simply explain the phenomenon of thermal expansion using the previous result. If the temperature increases, the average kinetic energy also increases, which means these particles tend to move faster. Their chemical bonds tend to loosen or break, which causes the length or volume of matter to increase slightly.
\begin{equation*}
	\overline{v^{2}}=\frac{3kT }{m}
\end{equation*}
We are also interested in the \textbf{average of the squares of the speeds} (or root-mean square, rms for short):
\begin{equation*}
	v_{\text{rms}}=\sqrt{\overline{v^{2}}}=\sqrt{\frac{3kT }{m}}
\end{equation*}
As you can see here, we have just done a bit of Statistical Physics. Because we do not know the \textbf{exact} state of each particle in a set of $10^{23}$ particles, we can still calculate their average velocity or energy\footnote[1]{Potential energy $U=mgz$ can be ignored, since $K>>U$.}. In the spirit of Statistical Physics, we solve the whole big problem of the enormous number of observations, without knowing the behavior or each individual particle.
% section The Kinetic Theory of Gases (end)
\section{The Equipartition Theorem}\label{sec:the_equipartition_theorem} % (fold)
The average translational kinetic energy is just a special case of a much more general result. In the previous discussion, we only considered the motion of particles moving in three-dimensional space. There are $3$ types of motion, directly related to $3$ types of kinetic energy, and each of them can be called a \textbf{degree of freedom}:
\begin{enumerate}
	\item Moving along $x$-axis: $ \overline{K}_{x}=\displaystyle\frac{1 }{2}m\overline{v_{x}^{2}}=\frac{kT }{2}$
	\item Moving along $y$-axis: $ \overline{K}_{y}=\displaystyle\frac{1 }{2}m\overline{v_{y}^{2}}=\frac{kT }{2}$
	\item Moving along $z$-axis: $ \overline{K}_{z}=\displaystyle\frac{1 }{2}m\overline{v_{z}^{2}}=\frac{kT }{2}$
\end{enumerate}
Other degrees of freedom include rotational motion and elastic potential energy (as stored in a structure like a spring), and what they have in common is that their energy functions are all \textbf{quadratic form}.
\begin{theorem}[The Equipartition Theorem]
	\label{The Equipartition Theorem}
	At temperature $T$, the average energy of any quadratic degree of freedom is $ \displaystyle\frac{1 }{2}kT $.
\end{theorem}
\begin{corollary}[Total Internal Energy]
	\label{Total Internal Energy}
	If a system contains $N$ particles, each with $f$ degrees of freedom, and the \textbf{total internal energy}\footnote[1]{Actually, there is a slight difference between total thermal energy and total internal energy; but I neglect them and keep it simple as possible.} is: $$ U=N.f. \frac{1 }{2}kT $$
\end{corollary}
The final problem is determining the value of $f$. For example, if our gas is monatomic, then we can easily conclude that  $f=3$. In the case of diatomic, $f=5$ because there are $5$ possible ways of motion:
\begin{enumerate}
	\item Moving along $x$-axis: $ \overline{K}_{x}=\displaystyle\frac{1 }{2}m\overline{v_{x}^{2}}=\frac{kT }{2}$
	\item Moving along $y$-axis: $ \overline{K}_{y}=\displaystyle\frac{1 }{2}m\overline{v_{y}^{2}}=\frac{kT }{2}$
	\item Moving along $z$-axis: $ \overline{K}_{z}=\displaystyle\frac{1 }{2}m\overline{v_{z}^{2}}=\frac{kT }{2}$
	\item Rotating around $x$-axis: $\overline{K}_{x,\text{ rotate}}=\displaystyle\frac{1}{2}I\omega_{x}^{2}=\frac{kT}{2}$
	\item Rotating around $y$-axis: $\overline{K}_{y,\text{ rotate}}=\displaystyle\frac{1}{2}I\omega_{y}^{2}=\frac{kT}{2}$
\end{enumerate}
Most polyatomic molecules can rotate about all three axes, so if our gas is polyatomic, then $f=6$. In the case of solid, $f=6$ too.
\begin{figure}[H]
	\begin{center}
		\includegraphics[width=0.8\textwidth]{~/PhyAI/chapter1/1_6.png}
	\end{center}
	\caption{Degrees of freedom illustration}\label{fig:1.6}
\end{figure}
Counting degrees of freedom of liquid is more complicated than gas or solid, since the intermolecular potential energies are not quadratic functions.
\begin{exmp}
	Calculate the internal energy of a bottle filled with $\ce{H2}$ at normal room temperature.
\end{exmp}
Assuming the normal room condition are $P=1$ (atm), $T_{C}=25^{\circ}$C and the bottle's volume is $V=500$ (ml). Using the idea gas law, we have:
$$ PV=NkT\Rightarrow N=\frac{PV }{kT} $$
Applying Corollary ~\ref{Total Internal Energy}, we obtain:
$$ U=N.f.\frac{kT }{2}=\frac{PV }{kT}.f.\frac{kT }{2}=\frac{PVf }{2}=\frac{101325.0.5.10^{-3}.5}{2}=126.65\text{ (J)} $$
% section The Equipartition Theorem (end)
\section{The First Law of Thermodynamics}\label{sec:the_first_law_of_thermodynamics} % (fold)
\subsection{Heat and Work}\label{Heat and Work}
Consider a piston filled with ideal gas\footnote[1]{Actually, you do not need to worry too much about the word "ideal", because most of the gases we encounter in real life behave similar to ideal gas model.}, now you heat it up with a candle, what will happen?
\begin{figure}[H]
	\begin{center}
		\includegraphics[width=0.7\textwidth]{~/PhyAI/chapter1/1_7.png}
	\end{center}
	\caption{Heat $Q$ is converting to work $W$ to push the piston away}\label{fig:1.7}
\end{figure}
I think we are all familiar with the definition of work in Classical Mechanics course: $$ W=\int\overrightarrow{F}d\overrightarrow{r} \text{ (J)}$$
Both work $W$ and heat $Q$ are $2$ types \textbf{energy}; but they have a slight but subtle difference:
\begin{itemize}
	\item Heat is defined as any \textbf{spontaneously flow of energy} from one object to another. As we discussed in the previous chapter, heat tends to flow from the higher temperature object to the lower, but \textbf{not the reverse direction}. We have not known the reason yet.
	\item Work is defined as any \textbf{other transfer of energy into or out of a system} (not spontaneous).
\end{itemize}
Notice that both heat and work refer to energy in \textbf{transmit}. It is meaningless to ask how much heat or work inside the system. We can only talk about how much heat or work has entered or leaf the system. According to the law of conservation of energy, energy can only be transformed or transferred from one to another. If we use the symbols $Q$ and $W$ to represent the amounts of energy that \textbf{enter} a system as heat and work; $U$ as the internal energy of system, then we can confidently state that:
\begin{center}\textit{External energy will increase the internal energy.}\end{center} Or written as an equation, this statement is: $$ \Delta U=Q+W\footnote[1]{Many other textbooks, including Halliday, define $W$ as work that \textbf{left} a system, so the sign of $W$ is negative: $$ \Delta U=Q-W $$ This convention works very great when solving heat engine or heat cycle problems.} $$
\begin{theorem}[The First Law of Thermodynamics]
	The increment in the internal energy is \textbf{equal} to the difference between heat accumulated by the system and the thermodynamics work done by it. $$ \Delta U=Q-W \footnote[2]{I read this statement from \href{https://en.wikipedia.org/wiki/First_law_of_thermodynamics}{Wikipedia}, and I think this is the best explanation for the first law.}$$
\end{theorem}
In this scope of this book, we will stick to the convention that I have mentioned above, to keep consistent, that means $W$ is representing the amount of energy that \textbf{enter} the system.
% section The First Law of Thermodynamics (end)
\subsection{Thermodynamics Process}\label{sec:thermodynamics_process} % (fold)

% section Thermodynamics Process (end)

% section Heat Cycle (end)
% chapter The First Law of Thermodynamics (end)
\end{document}
